\section{Balanced Mersenne multiplier localization}
\label{sec:mersenne-balanced-multiplier-depth}

We refine Section~\ref{sec:mersenne-multiplier-index-two-arm} in two
directions.  The one-copy multiplier cutoff can be moved closer to its
critical harmonic scale, and a published pointwise $p$-adic logarithm bound
removes a growing low-multiplier interval from the deep-valuation arm.  No
density conjecture is used.

For an odd prime $p$, retain
\[
 d_p=\ord_p(2),\qquad w_p=v_p(2^{d_p}-1),\qquad
 r_p=\frac{p-1}{d_p},
\]
and let $a_d=\sum_{d_p=d}(w_p-1)\log p$.  Fix a positive integer $k$ and a
real number $\eta>0$.  Put
\begin{equation}\label{eq:balanced-multiplier-scales}
 L_m=\log\log(3m),\qquad
 F_{m,\eta}=\log(3m)L_m^{1+\eta},\qquad
 H_{m,\eta}=\left\lfloor
 \sqrt{\frac{\log(3m)}{L_m^{1+\eta}}}\right\rfloor.
\end{equation}
Split the one-copy repeated support into
\begin{align*}
 U_d^{(\eta)}(m)
 &=\sum_{\substack{d_p=d,\ w_p\ge2\\
          p\le d^2/F_{m,\eta}}}\log p,\\
 B_d^{(\eta)}(m)
 &=\sum_{\substack{d_p=d,\ w_p\ge2\\
          p>d^2/F_{m,\eta}}}\log p,
\end{align*}
and split the deep excess by multiplier:
\begin{align*}
 V_d^{(\eta)}(m)
 &=\sum_{\substack{d_p=d,\ w_p\ge3\\
          r_p<H_{m,\eta}}}(w_p-2)\log p,\\
 G_d^{(\eta)}(m)
 &=\sum_{\substack{d_p=d,\ w_p\ge3\\
          r_p\ge H_{m,\eta}}}(w_p-2)\log p.
\end{align*}
The four terms give the exact identity
\begin{equation}\label{eq:balanced-multiplier-four-part}
 a_d=U_d^{(\eta)}(m)+B_d^{(\eta)}(m)
      +V_d^{(\eta)}(m)+G_d^{(\eta)}(m).
\end{equation}

\begin{proposition}[Adaptive one-copy cutoff]
\label{prop:balanced-adaptive-one-copy}
For fixed $k\ge1$ and $\eta>0$, eventually in $m$,
\[
 \sum_{\substack{d\mid m\\\log(m/d)<kL_m}}
 U_d^{(\eta)}(m)
 \le
 \frac{2m\log m}{F_{m,\eta}}(1+kL_m)
 =O_{k,\eta}\!\left(\frac{m}{L_m^\eta}\right)=o(m).
\]
\end{proposition}

\begin{proof}
Write $q=m/d$.  If $p$ occurs in $U_d^{(\eta)}(m)$, multiplier
injectivity and $p=1+d r_p$ give
\[
 r_p<\frac p d\le\frac d{F_{m,\eta}}.
\]
There are therefore at most $d/F_{m,\eta}$ such primes.  Eventually
$F_{m,\eta}\ge1$, so $p\le d^2\le m^2$ and
\[
 U_d^{(\eta)}(m)\le\frac{2d\log m}{F_{m,\eta}}.
\]
The window is $q<(\log(3m))^k$.  Hence
\[
 \sum U_d^{(\eta)}(m)
 \le\frac{2m\log m}{F_{m,\eta}}
    \sum_{\substack{q\mid m\\q<(\log(3m))^k}}\frac1q
 \le\frac{2m\log m}{F_{m,\eta}}(1+kL_m).
\]
Dividing by $m$ gives
$O_k((1+L_m)/L_m^{1+\eta})=O_k(L_m^{-\eta})$.
\end{proof}

Thus the exponent two on $L_m$ in the previous cutoff can be replaced by
$1+\eta$ for every fixed positive $\eta$.  The endpoint $\eta=0$ is not
claimed: the same calculation then gives only a bounded normalized error.
Relative to the old cutoff this is strictly sharper for $0<\eta<1$, equal at
$\eta=1$, and a valid but weaker localization for $\eta>1$.
That limitation is not an arithmetic counterexample, so the critical cutoff
remains an active route.

The deep estimate uses Yamada's unconditional theorem
\cite{Yamada2010}.  With
\[
                       C_Y=283\log3\log6,
\]
his Theorem~1.2, equation~(7), states for every prime $p$ that
\begin{equation}\label{eq:yamada-base-two}
 v_p(2^{p-1}-1)
 \le\left\lfloor C_Y\frac{p-1}{(\log p)^2}\right\rfloor+4.
\end{equation}

\begin{lemma}[Exact-order Yamada envelope]
\label{lem:exact-order-yamada-envelope}
For an odd prime $p$, with $d=d_p$, $r=r_p$, and $w=w_p\ge2$,
\[
 (w-2)\log p\le C_Y\frac{dr}{\log p}+2\log p.
\]
If $d>1$ and $r<H$, then
\begin{equation}\label{eq:yamada-multiplier-envelope}
 (w-2)\log p
 \le C_Y\frac{dr}{\log d}+2\log(1+dH).
\end{equation}
\end{lemma}

\begin{proof}
Since $p-1=dr$ and $0<r<p$, one has $p\nmid r$.  LTE gives
\[
 v_p(2^{p-1}-1)=v_p((2^d)^r-1)
 =v_p(2^d-1)+v_p(r)=w.
\]
Substitute this equality into \eqref{eq:yamada-base-two}, discard the floor,
subtract two, and multiply by $\log p$.  The final assertion uses
$p>d$ and $p=1+dr<1+dH$.
\end{proof}

\begin{lemma}[Triangular multiplier energy]
\label{lem:triangular-multiplier-energy}
If positive integral labels $r(x)<H$ are injective on a finite set $S$, then
\[
 |S|\le H-1,
 \qquad
 \sum_{x\in S}r(x)\le\frac{H(H-1)}2.
\]
Consequently, for $d>1$ and $H\ge2$,
\begin{equation}\label{eq:finite-deep-multiplier-bound}
 \sum_{\substack{d_p=d,\ w_p\ge3\\r_p<H}}(w_p-2)\log p
 \le
 \frac{C_Yd}{\log d}\frac{H(H-1)}2
 +2(H-1)\log(1+dH).
\end{equation}
\end{lemma}

\begin{proof}
The image of the labels lies in $\{1,\ldots,H-1\}$, proving both finite
bounds.  Apply \eqref{eq:yamada-multiplier-envelope} prime by prime and use
injectivity of $p\mapsto r_p$ in the exact-order fibre.
\end{proof}

\begin{theorem}[Low-multiplier deep mass is negligible]
\label{thm:balanced-low-deep-negligible}
For every fixed $k\ge1$ and $\eta>0$,
\[
 \sum_{\substack{d\mid m\\\log(m/d)<kL_m}}
       V_d^{(\eta)}(m)=o(m).
\]
More precisely, if $Q_m=(\log(3m))^k$ and $H=H_{m,\eta}$, then eventually
\begin{align}
 \sum V_d^{(\eta)}(m)
 \le{}&
 \frac{C_YmH(H-1)}{2\log(m/Q_m)}(1+kL_m)
 \notag\\
 &+2HQ_m\log(1+mH).
\label{eq:window-low-deep-bound}
\end{align}
\end{theorem}

\begin{proof}
In the window $q=m/d<Q_m$, hence $d>m/Q_m$.  Sum
\eqref{eq:finite-deep-multiplier-bound}.  The coefficient of the triangular
term is bounded using
\[
 \sum_{\substack{q\mid m\\q<Q_m}}d
 =m\sum_{\substack{q\mid m\\q<Q_m}}\frac1q
 \le m(1+kL_m),
\]
and $1/\log d\le1/\log(m/Q_m)$.  There are fewer than $Q_m$ integral
co-divisors, giving the second line of
\eqref{eq:window-low-deep-bound}.  Finally
\[
 H^2\le\frac{\log(3m)}{L_m^{1+\eta}},
 \qquad \log(m/Q_m)=\log m-kL_m\sim\log m.
\]
The first line divided by $m$ is $O_k(L_m^{-\eta})$, and the second line is
polylogarithmic in $m$, hence $o(m)$.
\end{proof}

The same slack $L_m^\eta$ therefore balances the normalized losses in
Proposition~\ref{prop:balanced-adaptive-one-copy} and
Theorem~\ref{thm:balanced-low-deep-negligible}.

\begin{theorem}[Balanced two-arm obstruction]
\label{thm:balanced-high-multiplier-two-arm}
Suppose the Mersenne endpoint $\log W_m=o(m)$ fails, and fix $\eta>0$.  Then
there are fixed $k\ge1$, $\epsilon>0$, and an unbounded sequence $m_j$ such
that for every fixed $0<\gamma<1/2$, after passage to a subsequence one same
alternative holds for all large $j$:
\begin{enumerate}
\item
\[
 \sum_{\substack{d\mid m_j\\\log(m_j/d)<kL_{m_j}}}
 G_d^{(\eta)}(m_j)\ge\gamma\epsilon m_j,
\]
and every counted prime has
$r_p\ge H_{m_j,\eta}$ and $p\ge1+d_pH_{m_j,\eta}$;
\item
\[
 \sum_{\substack{d\mid m_j\\\log(m_j/d)<kL_{m_j}}}
 B_d^{(\eta)}(m_j)\ge\gamma\epsilon m_j,
\]
and every counted prime has
\[
 d_p<\sqrt{p\log(3m_j)L_{m_j}^{1+\eta}}.
\]
\end{enumerate}
\end{theorem}

\begin{proof}
The fixed-window equivalence supplies a localized $a_d$-mass at least
$\epsilon m_j$.  Sum \eqref{eq:balanced-multiplier-four-part} and use the two
little-oh estimates above.  The $G^{(\eta)}$ and $B^{(\eta)}$ sums total
$(\epsilon-o(1))m_j$, so one carries every fixed share below one half on a
subsequence.  The displayed prime restrictions are respectively the
definition of $G^{(\eta)}$ and the rearrangement
$p>d_p^2/F_{m_j,\eta}$.
\end{proof}

Thus it now suffices to prove that the two surviving localized sums are
$o(m)$ for one fixed $\eta>0$ and every fixed $k$.  The first concerns deep
base-two lifts only after their order multiplier tends to infinity; the
second lies within a factor
$\sqrt{\log(3m)L_m^{1+\eta}}$ of the square-root order threshold.  Neither
estimate follows from the global Erd\H{o}s--Murty theorem
\cite{ErdosMurty1999} or the Murty--S\'eguin valuation dictionary
\cite{MurtySeguin2019}.

The exact scope of newer literature is also important.  Li--Zhao
\cite{LiZhao2026} fix one prime ideal and prove eventual stabilization of
successive prime-power order growth; their threshold depends on that prime
ideal and therefore does not bound initial depths uniformly as $p$ varies.
Large-sieve results for Fermat quotients such as Shparlinski
\cite{ShparlinskiFermatQuotients2011} average over intervals of the base
variable and do not specialize to fixed base two with the simultaneous
exact-order, divisor-window, and square-divisibility restrictions.  The
conditional non-Wieferich conclusions of Fellini--Murty
\cite{FelliniMurty2026} likewise do not close either surviving arm.

There are two exact sharpness boundaries.  For every integer $H\ge1$, the
complete positive packet $S_H=\{1,\ldots,H-1\}$ satisfies
\[
 \sum_{r\in S_H}r=\frac{H(H-1)}2.
\]
For any fixed $C$, choosing an integer $H>2C+1$ gives
$\sum_{r\in S_H}r>CH$; hence injectivity plus a pointwise linear envelope
does not imply a uniform $O(H)$ energy bound.  The member $H=4$ already
refutes the coefficient-one inequality.  At the two-arm level,
$(R,U,V,G,B)=(2,0,0,1,1)$ defeats every coefficient greater than one half.
Both are complete abstract counterexamples to the respective stronger
finite inferences.  In actual arithmetic,
\[
 \ord_{3511}(2)=1755,\qquad
 v_{3511}(2^{1755}-1)=2,\qquad
 \frac{3511-1}{1755}=2,
\]
so the universal assertion that every repeated exact-order prime has
multiplier at least three is also false with all premises present.  This
finite depth-two row does not refute either asymptotic surviving arm.

The companion Lean module formalizes triangular multiplier energy, the
positive-label capacity, the exact LTE transport from $p-1$ to the canonical
exact-order valuation, the affine weighted envelope and its realization on
the actual exact-order fibre, the two-controlled/two-surviving finite ledger,
the complete triangular sharpness family and the two-arm sharpness example,
and the full $3511$ order,
square-divisibility, and multiplier certificate.  Yamada's theorem and the
real asymptotic summation above remain cited paper mathematics; no valuation,
density, or abc assertion is introduced as a Lean axiom.
